\section{Introducción}

\subsection{Planteamiento del problema}
Los materiales compuestos reforzados con fibras son cada vez más utilizados en sectores de alto impacto como la aeronáutica, el transporte, la energía y la construcción, debido a su elevada relación resistencia–peso y a su versatilidad de diseño. Sin embargo, la eficiencia mecánica de estos materiales depende en gran medida de la calidad de la interfaz entre la fibra y la matriz, que es la encargada de transferir los esfuerzos. Una interfaz deficiente puede conducir a fallas prematuras, pérdida de rigidez y reducción de la vida útil de los componentes, lo que limita la competitividad de los compuestos en aplicaciones críticas \cite{HullClyne1996,KellyZweben2000}.

Para evaluar esta adhesión interfacial, se han desarrollado diversos métodos experimentales a escala micromecánica, entre los que destacan el \textit{Single-Fiber Pull-Out Test}, el \textit{Push-Out Test} y el \textit{Microdroplet Test}. Estos ensayos han sido ampliamente documentados y normalizados, permitiendo obtener estimaciones de la resistencia al corte interfacial (\textit{Interfacial Shear Strength}, IFSS). No obstante, su aplicación sobre fibras individuales presenta limitaciones importantes: alta dispersión de resultados, poca representatividad frente al comportamiento real de haces de fibras y dificultad para escalar a configuraciones prácticas de ingeniería \cite{DINSPEC19289}.

En este contexto, el \textit{Fiber-Bundle Pull-Out Test} (FBPO) se presenta como una alternativa experimental de gran potencial. Este ensayo consiste en extraer haces de fibras embebidos en una matriz, lo que permite evaluar la interacción de la interfaz a una escala más cercana a la real y obtener resultados con mayor robustez estadística. De León et al.\ \cite{DeLeon2022} propusieron un método de fabricación escalable para probetas de FBPO, utilizando moldes multicavidad y dispositivos de alineación geométrica, lo que garantiza control en la longitud de embebido, mejora la repetibilidad y permite la producción de grandes lotes de muestras. Dicho avance abre la posibilidad de implementar este ensayo en laboratorios convencionales y de generar datos sistemáticos para estudios comparativos de adhesión.

Asimismo, la implementación del FBPO encuentra respaldo en normativas afines que, aunque desarrolladas para fibras individuales o hilos textiles, entregan criterios transferibles. La \textbf{DIN SPEC 19289} define un procedimiento estandarizado para el \textit{single-fiber pull-out}, con énfasis en el control geométrico y en la alineación del ensayo. La \textbf{ISO 3341} establece lineamientos para la tracción de hilos de vidrio, incorporando prácticas de sujeción y acondicionamiento que pueden extrapolarse a haces de fibras. Por su parte, la \textbf{DIN 29965} entrega especificaciones para hilos de carbono en aplicaciones aeroespaciales, relevantes para la trazabilidad y caracterización del material utilizado en las probetas \cite{ISO3341,DIN29965}.

La pertinencia de abordar este problema radica en que la implementación del \textit{Fiber-Bundle Pull-Out Test} permitiría cerrar la brecha entre los ensayos micromecánicos tradicionales y las necesidades de caracterización a mesoescala, aportando a la mejora del diseño y la confiabilidad de materiales compuestos. Desde un punto de vista tecnológico, este avance permitirá contar con datos más representativos de la interacción fibra–matriz. En el ámbito económico, contribuirá a optimizar el uso de fibras de alto costo al reducir la incertidumbre en la etapa de diseño. En lo ambiental, una caracterización más precisa permitirá aumentar la vida útil de los componentes y reducir el desperdicio de materiales. Finalmente, en lo social y ético, fortalecer la confiabilidad de los materiales compuestos en aplicaciones críticas tiene un impacto directo en la seguridad y en la sostenibilidad de soluciones tecnológicas.

\subsection{Alcance}
El presente trabajo se enmarca en la implementación experimental del ensayo \textit{Fiber-Bundle Pull-Out} (FBPO) en materiales compuestos, con el objetivo de caracterizar la adhesión fibra–matriz a nivel de mesoescala. El alcance de la investigación considera el diseño y fabricación de mordazas, moldes y probetas mediante modelado CAD y técnicas de manufactura, la ejecución de ensayos bajo condiciones controladas en máquina de tracción, y el análisis estadístico de los resultados utilizando ANOVA. De esta forma, se busca establecer un protocolo experimental reproducible que permita evaluar de manera representativa la resistencia al corte interfacial.

Este estudio se limita al análisis experimental en laboratorio y no contempla el desarrollo de modelos numéricos, la evaluación de componentes estructurales a escala industrial ni el estudio de efectos de degradación ambiental como fatiga o envejecimiento. El horizonte temporal de ejecución corresponde a un período de cuatro meses, en el cual se espera generar un procedimiento validado y transferible que complemente las metodologías convencionales de caracterización interfacial.

\subsection{Objetivos}

\subsubsection{Objetivo General}
Implementar el ensayo \textit{Fiber-Bundle Pull-Out} en materiales compuestos, mediante el diseño y 
fabricación de un sistema experimental y la aplicación de una metodología adaptada a mesoescala, 
con el fin de determinar la resistencia al corte interfacial (IFSS) y validar su utilidad como técnica 
de caracterización confiable en el estudio de la adhesión fibra–matriz.

\subsubsection{Objetivos Específicos}
\begin{itemize}
    \item \textbf{OE1:} Diseñar y fabricar el sistema experimental a mesoescala —incluyendo mordazas, moldes y probetas— mediante modelado CAD, selección de materiales y técnicas de manufactura, con el propósito de asegurar precisión geométrica y reproducibilidad en los ensayos.
    \item \textbf{OE2:} Implementar y ejecutar el ensayo \textit{Fiber-Bundle Pull-Out} en probetas de material compuesto, utilizando equipos de tracción instrumentados bajo condiciones controladas, con el fin de obtener parámetros experimentales que describan la interacción fibra–matriz.
    \item \textbf{OE3:} Analizar y comparar los resultados experimentales mediante la aplicación de ANOVA y su contraste con antecedentes bibliográficos, para evaluar la resistencia al corte interfacial y validar la aplicabilidad del método como técnica de caracterización a mesoescala.
\end{itemize}


\subsection{Propuesta metodológica}

La metodología propuesta se estructura en tres fases principales, coherentes con los objetivos planteados. En primer lugar, se desarrollará el diseño y fabricación del sistema experimental a mesoescala, que comprende mordazas, moldes y probetas. Para ello se empleará modelado asistido por computadora (CAD), selección de materiales adecuados y técnicas de manufactura como impresión 3D y mecanizado, con el fin de garantizar precisión geométrica, control de la longitud de embebido y reproducibilidad de los ensayos.

En una segunda etapa se llevará a cabo la implementación del ensayo \textit{Fiber-Bundle Pull-Out}, utilizando las probetas fabricadas y un equipo de tracción instrumentado. Los ensayos se realizarán bajo condiciones controladas de carga y desplazamiento, registrando las curvas fuerza–desplazamiento y documentando los modos de falla observados en la interfaz. Esta fase permitirá obtener los parámetros experimentales necesarios para estimar la resistencia al corte interfacial (IFSS).

Finalmente, en la tercera etapa se procesarán y analizarán los resultados mediante técnicas estadísticas, en particular ANOVA, con el objetivo de identificar diferencias significativas entre condiciones de ensayo y contrastar los hallazgos con antecedentes bibliográficos. Con ello se busca validar la aplicabilidad del método a mesoescala y establecer un protocolo experimental confiable que complemente las metodologías tradicionales de caracterización interfacial.

